\newpage
% \section*{ВВЕДЕНИЕ}
\begin{center}
  \textbf{ВВЕДЕНИЕ}
\end{center}
\addcontentsline{toc}{section}{ВВЕДЕНИЕ}

Морская береговая линия, являясь местом совокупного взаимодействия ансамбля гео-, гидро- и аэромеханических процессов, представляет собой сложную и динамическую систему, обладающую, в результате влияния внешних сил, высокой степенью потенциального риска деформаций прибрежных территорий.
Указанные деформации, в свою очередь, влекут за собой утрату конструкционной целостности находящихся в прибрежных зонах инженерных и природных объектов, что, помимо экономического ущерба, провоцирует опасность и для человека.

При этом естественная тяга человека к водным объектам обуславливает заполнение прибрежных территорий наравне с объектами промышленного и инфраструктурного назначения рекреационными и жилыми зонами.
Распределение проявления негативных явлений во времени, в совокупности с неравномерностью распределения зон повышенного риска вдоль береговой линии, создают у внешнего наблюдателя иллюзию естественности и неизбежности этих процессов, выражающуюся в молчаливом принятии тезиса \enquote{вода камень точит}.
Лингвистическим антагонистом этого утверждения является известное \enquote{под лежачий камень вода не течет}, что формирует дугу диалектического противоречия, внести свой вклад в разрешение которого преследует, в идейном плане, цель данной работы.

Рассматривая практическую плоскость вопроса, следует опираться на тот факт, что основная масса удобных для своего освоения участков береговой линии была занята к настоящему моменту в ходе эволюционного заполнения прибрежных земель городским пространством.
Оставшиеся свободными для развития прибрежные территории требуют выполнения мероприятий по оценке своего потенциала.
Для обеспечения экономической эффективности освоения прибрежных территорий одним из факторов для такой оценки должен являться риск проявления негативных явлений, связанных с гидромеханическими процессами, накладывающими на землепользователя дополнительные траты на обустройство берегозащитных конструкций или ликвидацию последствий в случае пренебрежения этими мероприятиями.
В любом случае непреклонные силы природы возьмут свое, вопрос лишь в том, какую степень риска готовы принять ответственные землепользователь и государство, использующие свои ресурсы к достижению коллективной пользы.

Опираясь на сказанное, разработка общедоступного метода оценки прибрежных морских территорий по степени риска проявления негативных деформационных и гидромеханических процессов является актуальной задачей, решение которой повысит как экономическую эффективность мероприятий по их освоению, так и безопасность и сохранение названных земель при использовании.

\vspace{1em}
\textbf{Формулировка проблемы}

В градостроительных отношениях применяются методы формирования морских прибрежных территорий, которые не учитывают гидромеханических и геомеханических процессов вдоль морской береговой линии, что может привести к созданию потенциально опасной ситуации для жителей города и повлечь за собой негативные экономические эффекты.

\vspace{1em}
\textbf{Объект исследования}

Система морских прибрежных территорий, планируемая к градостроительному освоению.

\vspace{1em}
\textbf{Предмет исследования}

Метод оценки морских прибрежных территорий по степени риска проявлений негативных гидромеханических процессов.

\vspace{1em}
\textbf{Цель исследования}

Разработка и внедрение метода оценки морских прибрежных территорий по степени риска проявлений негативных гидромеханических процессов.

\vspace{1em}
\textbf{Идея работы}

Оценка риска негативных гидромеханических процессов за счет анализа волновых нагрузок на береговую линию.

\vspace{1em}
\textbf{Задачи исследования}

\begin{enumerate}

    \item Разработка метода оценки прибрежных территорий по уровню негативного воздействия, вызванного морским волнением.
    \item Анализ доступности и качества необходимых исходных данных, а так же разработка инструментов доступа и работы с ними.
    \item Алгоритмическая и программная реализация разработанного метода.
    \item Анализ результатов применения метода для локальной задачи.

\end{enumerate}

\vspace{1em}
\textbf{Методология и методы}

\begin{enumerate}
    \item Методы научного поиска и анализа.
    \item Методы сбора и обработки геопространственных данных.
    \item Методы математической статистики.
\end{enumerate}

\vspace{1em}
\textbf{Новизна исследования}

Разработанные в ходе исследования методики получения геопространственных данных и вычисленные для морских прибрежных территорий индексы уязвимости позволят повысить качество научного описания состояния прибрежных территорий.

\vspace{1em}
\textbf{Теоретическая и практическая значимость}

\begin{enumerate}
\item Разработанные методы получения геопространственных данных могут быть использованы в смежных по тематике научных и инженерных работах.
\item Предложенные показатели волнового воздействия на берег могут служить основанием для совершенствования регламентов контроля и обслуживания городских берегоукрепительных и гидротехнических сооружений.
\item Распределение разработанных волновых индексов вдоль неурбанизированных морских прибрежных территорий может служить одним из критериев для выделения перспективных локаций для основания и развития территорий.
\end{enumerate}

\vspace{1em}
\textbf{Соответствие направлению подготовки}

Выполненный в ходе НИР анализ уязвимости морских береговых линий повысит эффективность территориального планирования, позволяя оптимизировать затраты на гидротехнические сооружения, время на подготовку проектных материалов, трудовые ресурсы, необходимые для их формирования, и сократить, в перспективе, траты на ликвидацию негативных гидромеханических событий, что согласуется с паспортом направления 27.04.07 «Наукоемкие технологии и экономика инноваций».
